Let \(g=4sr/(4sr+2s+r)\).  Since all denominators are positive,
\[
g<\beta_1
\quad\Longleftrightarrow\quad
4s(2r+1)<4sr+2s+r
\quad\Longleftrightarrow\quad
s<\frac{r}{4r+2}.
\]
Thus \(\gamma_1=g<\beta_1\) below the threshold and \(\gamma_1=\beta_1\) at or above it.  Moreover,
\[
g-\alpha_1
=\frac{s\{s(4r-2)+3r\}}{(4sr+2s+r)(2s+1)}.
\]
If \(r\geq1/2\), its numerator is positive.  If \(0<r<1/2\) and \(s<r/(4r+2)\), then
\[
s(2-4r)<\frac{r(2-4r)}{4r+2}<r<3r,
\]
so again \(s(4r-2)+3r>0\).  Hence \(\alpha_1\leq g=\gamma_1\) in the rough-mean regime.  Below the threshold both \(\alpha_1\wedge\gamma_1\) and \(\alpha_1\wedge\beta_1\) equal \(\alpha_1\); at or above it \(\gamma_1=\beta_1\).  This proves the claimed identity.